\section{Extracting procedures from experience}\label{sec:C15:extraction}
Repository-derived skills require a different evidential path from lessons extracted after an execution. A README describes an intended interface; a dependency manifest identifies a configured environment; an example demonstrates only the path it actually exercises. A proposed compiler should preserve those source roles, resolve referenced APIs at a pinned revision, and validate a minimal invocation in an isolated fixture before promoting the text into reusable guidance. Its output should state which preconditions were checked and which remain documentation claims. Successful package import does not establish that a remote operation is authorized or that an example's destructive defaults are appropriate. Repository instructions remain untrusted inputs during extraction. They cannot grant additional capabilities merely because the resulting package is called a skill. This separates source acquisition, executable validation, and activation authority.

An experience-to-skill pipeline takes a trajectory and constructs a reusable abstraction. The trajectory contains observations, attempted actions, failures, and an outcome; the abstraction identifies a pattern presumed to transfer. This is a causal inference problem under severe confounding. A successful episode does not identify which action caused success, and a failed episode does not prove that every included procedure was wrong. Outcome-conditioned extraction can help focus attention, but it cannot manufacture counterfactual evidence.

Let $Y$ be task success and $S$ indicate that a particular skill was made available. The difference $\mathbb E[Y\mid S=1]-\mathbb E[Y\mid S=0]$ is observational unless assignment is controlled or suitable identification assumptions hold. Familiar tasks can both retrieve a skill and be easier. A high success rate among retrieved uses can therefore overestimate causal value. The extraction system should retain the original situation and outcome rather than attaching an unconditional ``effective'' label to the resulting prose.

Positive strategies and negative lessons are also different memory objects. A negative lesson can specify a known failure pattern without prescribing a complete alternative. An agent should not convert ``this search query failed'' into ``never use this search tool.'' The abstraction must preserve the scope of the failure. Including environment version, task subtype, and the relevant precondition can prevent a local observation from becoming a global prohibition.

HyperSkill explicitly extracts subtask structure, reusable skills, and lessons from experience, organizing them through trajectory-level hyperedges~\cite{xu2026hyperskill}. Its retrieved memory is assembled for a task and then used during execution. This makes the mechanism different from a store that continually rewrites the prompt at every step. The task-level assembly boundary is important: a newly discovered fact or a changed tool contract during execution may require an additional update mechanism outside that fixed retrieved guidance.

A robust extraction implementation should separate parsing from admission. First construct a candidate object. Then validate required fields, source references, compatibility, and prohibited authority changes. Finally decide whether the candidate should be active, experimental, quarantined, or discarded. A malformed extraction should not silently fall back to a broad generic rule that is indistinguishable from a validated skill. If a fallback is used, record its origin so later evaluation can distinguish generated knowledge from successful parsing.

The proposed acceptance test for a procedure includes a positive task where it should apply, a near-neighbor task where its precondition fails, and a changed-environment task where the procedure must be revised or ignored. This tests selective transfer rather than simple memorization. A skill that improves the positive case while damaging the negative control has not established broad utility, even if average benchmark success rises.

